\documentclass[11pt]{article}
\usepackage[margin=1in]{geometry}
\usepackage{hyperref}
\title{Learning from Imbalanced Data}
\author{Devis W. Saputra}
\begin{document}
\maketitle

\begin{abstract}
I create a controlled class imbalance from real Wisconsin Diagnostic Breast Cancer observations and compare standard and class-weighted classifiers. The experiment focuses on Average Precision and F1 instead of raw accuracy.
\end{abstract}

\section{Question}
How much does class weighting change performance when one class is deliberately made much less common?

\section{Data and Method}
I keep all 212 class-0 observations and a seeded subset of 35 class-1 observations. I compare standard logistic regression, balanced logistic regression, and a balanced Random Forest with 350 trees.

\section{Results}
The logistic model reaches 0.9924 Average Precision and 0.9091 F1. Balanced logistic regression keeps the same Average Precision and improves F1 to 0.9565. The balanced Random Forest reaches 1.0000 for both metrics on this split.

\section{Limitations}
The class balance is created for the experiment and does not represent real medical prevalence. The held-out set is small, so the perfect Random Forest result should not be generalized.

\section{Reproduction}
Install \texttt{requirements.txt} and run \texttt{python src/run\_experiment.py}.

\bibliographystyle{plain}
\bibliography{references}
\end{document}
